Consider the case of a vortex within the fluid. Its vorticity satisfies
\[
 \partial_t\omega+(u\cdot\nabla)\omega
 =(\omega\cdot\nabla)u+\nu\Delta\omega,
 \qquad \omega=\nabla\times u.
\]

Strain within the fluid's flow can naturally stretch a vortex. For axial
strain rate \(\sigma>0\), tube length \(\ell\) grows as
\[
 \dot\ell=\sigma\ell>0.
\]
When that happens, incompressibility forces a material tube to narrow as it lengthens.
If circulation persists strongly enough through that narrowing, the fluid
circulates faster around a smaller core.
For a locally circular tube of radius \(r\) and circulation \(\Gamma\),
the radius and mean swirl speed \(v\) satisfy
\[
 r\propto\ell^{-1/2},
 \qquad v=\frac{|\Gamma|}{2\pi r}.
\]

For a material tube following a chosen group of fluid particles: a vortex
core is the region where rotation is concentrated. Stretching squeezes
those particles inward. Viscosity diffuses momentum between neighbouring
regions of fluid, reducing velocity differences and spreading that rotation
through the surrounding fluid.
The diffusion distance \(r_{\mathrm{diff}}\) and viscous contribution are
\[
 r_{\mathrm{diff}}\sim\sqrt{\nu t},
 \qquad
 \langle\omega,\nu\Delta\omega\rangle
 =-\nu\|\nabla\omega\|_2^2\le0.
\]
The material tube can narrow while some rotation spreads beyond its
boundary, reducing the circulation enclosed by it. So narrowing the vortex
core makes the fluid circulate faster only if the enclosed circulation does
not weaken too quickly.

